\capitulo{6}{Trabajos relacionados}

Este capítulo presenta una revisión de trabajos y herramientas relacionadas con la temática del proyecto. Sus finalidades son: 
\begin{itemize}
	\item Situar la propuesta desarrollada dentro del panorama existente de soluciones de gestión de reservas y recursos.
	\item Identificar qué necesidades ya se encuentran cubiertas por otras aplicaciones, y en qué aspectos la solución propuesta aporta un enfoque diferencial.
\end{itemize}

A diferencia de otros trabajos más orientados a la investigación teórica o el análisis algorítmico, en este caso, resulta importante comparar la aplicación desarrollada con herramientas ya existentes en el ámbito de la reserva de espacios, su gestión y la administración de recursos compartidos. Por ello, este apartado se centra en soluciones representativas que, aun perteneciendo a contextos parcialmente distintos, comparten cierta parte de las funcionalidades abordadas en el proyecto.

\section{Herramientas de reserva y gestión de citas}
\subsection{Bookitit}
Bookitit es una plataforma comercial de reservas \textit{online} orientada principalmente a la gestión de citas, agendas y servicios. Entre sus funcionalidades más destacadas se encuentran la reserva \textit{online} por parte de clientes, la gestión manual de citas introducidas presencialmente o por teléfono, la administración de calendarios, notificaciones y recordatorios, así como la sincronización con servicios externos como \textit{Google Calendar}.

Además, la plataforma contempla características adicionales como la gestión de clases o reservas grupales, personalización del \textit{widget} de reserva e integración con páginas web y redes sociales. Estas capacidades la convierten en una solución versátil para sectores como salud, estética, formación o actividades dirigidas.

No obstante, su orientación principal se sitúa en la gestión comercial de citas y servicios, más que en la organización académica de docencia reglada. Aunque alguna de sus funciones pueden resultar comparables a la gestión de reservas puntuales planteada en este proyecto, no está específicamente creada para trabajar con horarios docentes, calendarios académicos, reservas periódicas asociadas a asignaturas o restricciones académicas.

\section{Sistemas de reservas de recursos}
\subsection{LibreBooking}
LibreBooking es una solución de código abierto orientada a la reserva y planificación de recursos compartidos dentro de organizaciones. Su propósito principal es permitir la gestión de reservas de recursos mediante una interfaz flexible, extendible y adaptada a múltiples tipos de activos, como salas, equipamiento o recursos comunes.

Frente a herramientas más centradas en la cita previa comercial, LibreBooking se aproxima más al problema tratado en este proyecto, al trabajar directamente con la reserva de recursos y gestión de disponibilidad. Además, su carácter abierto y extensible hace que resulte especialmente interesante como referencia desde el punto de vista de la arquitectura y del enfoque de gestión.

Sin embargo, se trata de una herramienta generalista de reserva de recursos. Aunque puede resultar útil para reservar aulas o espacios, no incorpora de forma nativa aspectos específicos del dominio universitario como la conversión automática de horarios académicos en reservas periódicas, la diferenciación entre tipo de usuarios institucionales, la integración con autenticación universitaria, la gestión detallada de roles académicos o la gestión del calendario académico.

